\chapter{Supporting analyses}
\label{appendix:supporting-analyses}

\section{Prior-art search methodology for \S6.3.1}
\label{appendix:prior-art-search}

This appendix documents the prior-art scan that informs the patentability analysis in \Cref{sec:ent-patentability}. It is presented as a methodological appendix, not as a substitute for a formal freedom-to-operate (FTO) opinion by patent counsel.

\subsection*{Calibration disclaimer}

This is an academic prior-art scan, not a freedom-to-operate opinion or legal advice. The search corpus is publicly indexed Google Patents, Espacenet and USPTO Public Search; no fee-based databases (Derwent, PatBase, IFI Claims, IPlytics) were consulted. All conclusions are framed as \emph{appears}, \emph{based on this scan}, \emph{likely}. An issued patent identified below is not necessarily blocking: a real FTO must read individual claims. Conversely, a journal or press disclosure is prior art for novelty under \S102 / Art.~54 EPC but is not an infringement risk.

\subsection*{Search engines}

\begin{itemize}
    \item Google Patents (\texttt{patents.google.com}) --- full-text plus citation graph; both granted and published-application records.
    \item Espacenet (EPO/WIPO/UKIPO full-text), cross-checked via Google Patents.
    \item USPTO Patent Public Search --- confirmation of US grant numbers.
    \item Justia Patents --- assignee and classification index.
    \item Secondary press/aggregator sources (Patently Apple, nweon Patent, EurekAlert, GlobeNewswire) for assignee--product linkage.
\end{itemize}

\subsection*{Search strings (representative)}

\begin{itemize}
    \item \texttt{"beamsplitter" eye tracking pupillometry patent}
    \item \texttt{"dichroic" "850 nm" eye tracking head-mounted display patent}
    \item \texttt{"hot mirror" eye tracking patent VR HMD pupil tracking}
    \item \texttt{"longpass dichroic" beamsplitter eye imaging patent pupil 750 nm}
    \item \texttt{"binocular vergence" pupillometry concussion patent headset}
    \item \texttt{"pupillary light reflex" patent headset goggles concussion screening}
    \item Assignee queries: NeurOptics, Apple Inc., Magic Leap, Oculus/Meta/Facebook, Tobii, SyncThink/Neurosync, Oculogica, Neuro Kinetics/Neurolign/Spryson, 3M, Elbit Systems, Novutz LLC, Indiana University Research and Technology Corp, Purdue Research Foundation, University of Miami.
\end{itemize}

\subsection*{Classification codes used (CPC / IPC)}

\begin{itemize}
    \item A61B 3/113 --- Objective instruments for examining eye movement.
    \item A61B 3/14 --- Arrangements specially adapted for taking photographs of the eye.
    \item A61B 5/0059 --- Detecting/measuring for diagnostic purposes by tracking eye-movement pattern.
    \item G02B 27/14 --- Beam splitting or combining systems based on path interference.
    \item G02B 27/01 --- Head-up displays / head-mounted displays.
    \item G06V 40/19 --- Sensor that detects gaze direction.
\end{itemize}

\subsection*{Date range}

Searches were not date-bounded. Hits range from 1992 (US5,345,281A) to 2025 grants (e.g.\ US12,311,584). The planned priority filing window is 2026--2027; any patent or publication pre-dating priority is potential anticipating prior art.

\subsection*{Coverage limits}

No fee-based prior-art database (Derwent, PatBase, IFI Claims, IPlytics) was consulted. All queries are keyword-based; no structured semantic search. Non-patent literature (NPL) was scanned only opportunistically through patent-database back-references and through targeted academic searches for BrainEye, FOVE-VR, and concussion-screening publications. Foreign-language patents (CN, JP, KR) were reachable via Google Patents machine translation but were not exhaustively scanned. A real FTO would require reading independent claims of each issued patent element-by-element against the proposed device's draft claim set.

\subsection*{Most-relevant references identified}

The twelve most novelty-impacting records identified are catalogued, with assignee, filing/priority date, relevance, and novelty-impact assessment, in the internal audit working file (not reproduced in this report for length reasons; available on request). The five highest-impact entries are: \citet{taboada1994}, \citet{neuropticspatent2017}, \citet{neurokineticshmd2017}, \citet{krueger2014ocular}, and \citet{syncthinkpatent2015}; each is cited in \Cref{sec:ent-patentability} above.
